\section{Casos prácticos y ruta de implementación}

\subsection{Caso: sistema Multi-hop DSPy}

Un pipeline multi-hop típico primero recupera, después razona, luego verifica y después invoca herramientas. DSPy te permite abstraer cada uno de estos pasos por separado en un module, y luego usar un optimizer para ayudarte a buscar combinaciones de prompt y ejemplos. En la práctica te encontrarás con problemas de sincronización en la línea de tiempo entre `chain-of-thought` y `tool invocation`.

\subsection{Gobernanza y documentación}

Se recomienda agregar un documento `DSPy.md` nuevo en el repo, que incluya la versión de deployment, los enlaces del Action Timeline y del Evidence Matrix, así como la security checklist. En cada release, `DSPy.md` se incluye como parte de la release note y se entrega a QA para su revisión.

\begin{importantbox}{La documentación es gobernanza}
El código de DSPy tiene dos capas más que un pipeline normal: DSL + optimizer log. Si no se documenta, aparecerá entre los desarrolladores la fricción de ``no sé por qué el optimizer cambió el prompt''. La documentación debe incluir: signature, metric, condiciones de fallback, contacto del owner.
\end{importantbox}

\subsection{Resumen de la sección}

Descomponer los sistemas compuestos de gran escala en casos reutilizables, y mantener para ellos documentación en texto plano, es un paso clave para pasar de la prueba de concepto a la operación real.

